% ---------------------------------------------------------------------------
% Response to R1 Major 3 -- JEDC revision
% Compile with pdfLaTeX. No figures required; standard packages only.
% Set to fit two pages. To paste into an existing response document, take
% everything between the BODY BEGINS and BODY ENDS markers.
% ---------------------------------------------------------------------------
\documentclass[11pt]{article}

\usepackage[utf8]{inputenc}
\usepackage[T1]{fontenc}
\usepackage{lmodern}
\usepackage[margin=1in]{geometry}
\usepackage{booktabs}
\usepackage{amsmath}
\usepackage{microtype}
\usepackage[hidelinks]{hyperref}

\setlength{\parskip}{0.4em}
\setlength{\parindent}{0pt}

\title{\textbf{R1, Major 3}\\[0.3em]
       \large\normalfont\itshape Disciplining the substitution elasticity}
\author{}
\date{}

\begin{document}
\maketitle
\vspace{-3em}

% =========================== BODY BEGINS ===================================

\begin{quote}\itshape
Provide sensitivity analysis and, if possible, discipline the elasticity using
evidence on energy demand or imported-input substitution.
\end{quote}

We have done this, and it changed the calibration: the elasticity is now
estimated from our own import panel rather than assumed, and the estimate
rejects the value we submitted. The share parameter $\alpha$ moved with it, from
0.2 to 0.000058.

\textbf{Why $\alpha$ moved by so much.} $\alpha$ is not the oil share; its units
depend on the elasticity. Holding it fixed while changing the elasticity would
quietly change the oil share of GDP, one of our calibration targets, so each
elasticity needs its own $\alpha$. Across the six elasticities below, $\alpha$
varies by a factor of 1,800 while the simulated oil bill comes out at 10.08\% of
GDP in every case.

\section*{What the data say, and why we adopt 0.25}

The oil share identifies the elasticity directly: if oil is hard to substitute,
the oil bill takes a \emph{larger} share of GDP when oil is dear; if it is easy
to substitute, the share falls. In our panel of 34 oil importers, 1995--2015,
the share clearly rises with the price. Four specifications --- within country,
collapsed to a time series, collapsed with a time trend, and first differences
--- all give positive coefficients, implying elasticities of 0.39, 0.37, 0.04
and 0.16, which average 0.244. The baseline coefficient is $+0.606$ (s.e.\
0.083), putting the elasticity between 0.23 and 0.56. The elasticity of 2 we
submitted implies the opposite sign altogether --- oil bills \emph{falling} as a
share of GDP as oil gets dearer --- so both 2 and 1 are rejected: oil and the
domestic input are complements. We adopt \textbf{0.25} as a central value across
specifications rather than the most favourable one. (The sign flips if oil
exporters are left in the sample, so the importer definition matters, and we now
say so in the text.)

\section*{How the model responds across the range}

We solved the model at six elasticities, re-matching the discount factor and the
default-cost parameter each time so that every version still hits the spread and
debt targets, then simulated 750,000 quarters at each.

\begin{table}[h]
\centering
\footnotesize
\setlength{\tabcolsep}{5pt}
\begin{tabular}{ccccccc}
\toprule
elasticity & oil & value of & oil bill & output & default rate & mean \\
           & imports & oil imports & / GDP & & \% p.a. & spread, \% \\
\midrule
0.20          & 0.8159 & 0.9681 & 10.079 & 9.6654 & 3.481 (0.044) & 3.713 \\
\textbf{0.25} & \textbf{0.8173} & \textbf{0.9593} & \textbf{10.079}
              & \textbf{9.5817} & \textbf{3.345 (0.043)} & \textbf{3.589} \\
0.30          & 0.8177 & 0.9495 & 10.079 & 9.4872 & 3.469 (0.044) & 3.716 \\
0.35          & 0.8174 & 0.9391 & 10.080 & 9.3838 & 3.325 (0.043) & 3.664 \\
0.40          & 0.8159 & 0.9275 & 10.079 & 9.2687 & 3.224 (0.043) & 3.549 \\
0.50          & 0.8105 & 0.9017 & 10.079 & 9.0104 & 3.201 (0.043) & 3.567 \\
\bottomrule
\end{tabular}
\end{table}

The quantity of oil imported hardly moves --- under 1\% across the whole range.
Output is about 7\% lower at the top of the range than at the bottom, with the
oil bill falling in step; that is the production function being
re-parameterised at a fixed oil share, not a change in behaviour. The result
that matters is the last two columns: \textbf{the default rate and the spread
are flat}, 3.20\% to 3.48\% a year and 355 to 372 basis points, with no trend in
either and the largest gap from the adopted elasticity about two standard
errors. The paper's credit-risk results do not depend on where in this range the
elasticity is set.

The welfare number is not robust in the same way, and we now say so. Removing
all oil price risk is worth $+0.47\%$ of consumption at the adopted elasticity
but only $+0.08\%$ at 0.50, and it moves around in between, so a claim about its
size is as much a claim about the elasticity as about oil risk. The spread
result is the one we lead with: 181 to 224 basis points throughout.

An elasticity of 2 does worse than widen the error bars --- it reverses the
paper's own conclusion. A country that substitutes easily buys less oil when oil
is dear, so its oil bill is capped and price swings become a benefit rather than
a cost (Waugh, 1944). On our earlier solution grid the welfare gain turns
negative above 0.5 and reaches $-2.8\%$ of consumption at an elasticity of 2,
against the $+0.2\%$ to $+1.1\%$ the paper claims. The calibration we submitted
therefore contradicted our own headline result.

\section*{One elasticity for 34 different economies}

Estimating the share relationship country by country gives 34 separate
elasticities: median 0.38, interquartile range 0.19 to 0.65. The pooled estimate
is not the work of a few countries --- it sits close to the median and to the
mean weighted by oil import share (0.31). Thirty-two of the 34 are complements;
19 cannot reject 0.25, and only 3 cannot reject 2. Much of the visible spread is
measurement error rather than real difference between countries: the variance of
the 34 estimates is 0.349 against an average sampling variance of 0.219, leaving
genuine differences with a standard deviation near \textbf{0.36 rather than
0.59}.

What matters for the paper is whether the elasticity differs with oil
dependence, the dimension the mechanism runs through. It does not: the most
oil-dependent third of the sample averages 0.27 and the least 0.20, and
regressing the country estimates on oil dependence gives $-0.15$ with a standard
error of 0.17. Imposing one elasticity therefore does not systematically
misprice the economies that drive the result. These estimates are in any case
too imprecise to calibrate 34 economies separately --- five come out negative
--- so we report the dispersion rather than propose a heterogeneous version.

\bigskip
\hrule
\bigskip

{\footnotesize
Spreads in the table are averages over the whole simulation; elsewhere we
average over the quarters before a default, 3.59\% against 4.08\% here, so the
two are not comparable. Everything for this comment is in
\texttt{model/economies\_common\_check/}:
\texttt{JEDC\_elasticity.m}, which writes every table, and the subfolder
\texttt{elasticity\_sweep/} holding the scripts that solved the six economies,
the run record, and the tables themselves. The six Fortran
solves are not re-run; \texttt{sweep\_resp/responses.csv} is their record, and
the 0.25 row is rebuilt from
\texttt{model/economies\_common\_check/benchmark} and compared.
\par}

% ============================ BODY ENDS ====================================

\end{document}
